\chapter{What is artificial life?}

\section{A specimen}

You can find a video of one on YouTube without much trouble. A blob,
perhaps a few hundred pixels across, smudged at the edges, pulsing on
a thirty-frame-per-second loop. It moves. Sometimes it stretches, finds
something that looks like food, and engulfs it. Sometimes it splits in
two, and the daughter blobs drift apart and live their own lives.
Sometimes it shrinks, fades, and dies in a way that is genuinely sad
to watch. The whole thing runs on a laptop. The blob is not alive. The
system that produces it is a few hundred lines of code.

This is Lenia, and it is the kind of thing that gets people into
artificial life. The puzzle is not whether the blob is \emph{really}
alive --- it is not, in any sense that requires defending --- but
whether the things our brains do when we watch it are misfiring or
tracking something real. Pattern recognition? Kinship recognition?
Theory of mind? All of those, probably. But also: there is a structure
on the screen that maintains itself, responds to its surroundings,
reproduces, and dies. Those are not nothing. The question of artificial
life is how to read that \emph{something} without dismissing the
apparent aliveness as projection or asserting it as literal.

The field has been at this for fifty years. It has produced many
specimens. This chapter introduces five of them --- not as a survey,
but as a set of test cases. Each one will reappear later in the book as
a worked example, and each one isolates a different facet of what we
mean by \emph{emergent}, \emph{novel}, or \emph{alive}. Treat the
chapter as a tour through a small cabinet of curiosities. The
vocabulary for reading them properly comes in Chapter 4.

\section{Conway's Game of Life (1970)}

A two-dimensional grid of cells, each on or off. Each step, a cell with
two or three live neighbours stays on; an off cell with exactly three
live neighbours turns on; everything else goes off. That is the entire
rule.

Run it from a random initial state and watch. After a few hundred steps
the field is dotted with stable shapes --- a $2 \times 2$ block, a
``blinker'' that oscillates between horizontal and vertical --- and
travelling shapes --- a ``glider'' that moves diagonally one square
every four steps. Most cells flicker out into one of these endpoints.
There is a ``glider gun,'' a configuration that emits a glider every
thirty steps forever, and this turns out to be the ingredient you need
to prove that the Game of Life is Turing-complete. You can build a
calculator in it, given enough patience and grid space.

What is interesting about the Game of Life is not the calculator. That
was a payoff for showing the system is rich enough to express anything
computable, but the computational construction is brittle and not what
people \emph{do} with the system. What is interesting is the gliders.
Conway's rule was chosen by hand, after some experimentation, to
produce a system that was neither trivial --- everything dies,
everything stabilises into static blobs --- nor saturated --- the
field fills with noise. The gliders are durable, recognisable patterns
that no part of the rule explicitly references. They are
\emph{emergent} in the loosest sense of the word: a feature of the
system that is not anywhere in the system's specification.

We will come back to the Game of Life when we want to make this loose
sense of emergence precise. The gliders are also handy for asking:
what does it take for a pattern to count as a \emph{thing}? The book's
fifth slot --- the observer --- will turn out to be where that question
gets answered.

\section{Reynolds' Boids (1987)}

Each agent (``boid'') has a position and a velocity in a two- or
three-dimensional world. Each step, each boid adjusts its velocity
according to three rules: steer to match the average heading of nearby
boids (alignment), steer toward the average position of nearby boids
(cohesion), and steer away from boids that are too close (separation).
That is the entire algorithm.

Watch a population of fifty boids on a screen and you see flocking. Not
the rules --- the flocking. The flock has a shape. It turns coherently.
It splits to avoid an obstacle and re-forms on the other side. None of
that is in the rules; the rules say only what each individual boid does
to its neighbours.

Boids was Craig Reynolds' computer-graphics solution to the problem of
animating flocks of birds in 1987, and it has been used in nearly every
film involving a non-trivial flock since \textit{Batman Returns} in
1992. But it is also a paradigm case for a particular flavour of
emergence: collective behaviour that arises from local rules. The flock
is the thing the audience sees, and the flock is real in some sense ---
you can ask how fast it is moving, where its centre is, how cohesive it
is --- but the simulation contains no representation of it. There is no
flock object. There are only boids.

Boids will reappear in the book whenever we want to talk about the gap
between what a system does at the level of its individual components
and what it does as a whole. The book's fourth slot --- the interaction
topology --- will turn out to control how aggressively that gap opens.

\section{Lenia (Chan, 2019)}

Bert Chan's Lenia is a continuous generalisation of the Game of Life.
The grid is replaced by a continuous field, the binary cell-state is
replaced by a real number in the unit interval, and the discrete update
rule is replaced by a smooth function of a smoothly weighted
neighbourhood. The result is a system whose behaviour depends
sensitively on a handful of parameters: a kernel shape, a growth
function, a time step. Pick parameters in a small fraction of the
parameter space and you get the blob from the chapter's opening.

The blobs Lenia produces are called ``creatures'' in the literature,
and the scare quotes are doing real work. They have boundaries that
move. They have internal structure that persists as they translate
across the field. Two creatures from the same parameter set can collide
and either annihilate, merge, or push past each other depending on the
angle. Some parameter sets produce creatures that visibly reproduce
--- a parent creature divides into two daughters with the same
morphology --- although whether the reproduction is faithful enough to
support evolution is a question Lenia mostly leaves open.

Lenia is the book's running example. There are several reasons. It is
continuous, which lets us use the language of dynamical systems without
violence. Its parameter space is small enough that a reader can explore
it on a laptop and large enough to produce genuine surprises. And the
creatures it produces sit exactly at the seam between ``obviously a
pattern in a dynamical system'' and ``something that resembles an
organism in a way that makes you uncomfortable to dismiss.'' Every one
of the book's slots --- substrate, variation, viability, topology,
observer --- has a clean Lenia reading.

\section{Tierra (Ray, 1991)}

Tom Ray was an evolutionary biologist who, in the late 1980s, decided
to build a digital ecosystem. Tierra is a virtual machine --- a piece
of simulated computer hardware --- running self-replicating programs in
shared memory. Each program is a short sequence of instructions; the
simulation runs all programs in parallel, allocates CPU time, and lets
programs whose instructions copy themselves into a free region of
memory reproduce. Mutations are introduced by occasional bit-flips.
Programs that fail to replicate get culled when memory runs out.

Ray seeded Tierra with a single hand-written 80-instruction program,
hit go, and went to lunch. When he came back, the memory was full of
variants. Most were 80 instructions, slight mutations of the seed. But
there were also 45-instruction programs --- \emph{parasites} --- that
lacked the copying machinery themselves but used a pointer to find a
nearby full-length program and borrowed its copying machinery to
reproduce. There were 79-instruction programs that had evolved partial
defences against parasites. There were hyper-parasites that exploited
parasites.

For a few hours, Tierra was the most exciting thing in artificial life.
It seemed to be doing what Darwinian evolution was supposed to do:
producing genuine novelty, an arms race, ecological roles that no one
had specified. Then, after a while, it stopped. The system reached a
kind of dynamic equilibrium and stayed there. The novelty production
rate dropped to near zero. New variants kept appearing, but they were
variations within a settled ecology, not new ecological roles.

Tierra is the book's exemplar for a deep problem: many artificial-life
systems start out producing novelty, and then they stop.
\emph{Open-endedness} --- the indefinite production of genuinely new
structure --- is something natural evolution seems to do and that we do
not yet know how to engineer. The book's third slot, viability, and the
discussion of the open-endedness metric $\Omega$ in Chapter 8, will
both circle back to Tierra.

\section{Evolutionary computation (Holland, 1975)}

John Holland's \textit{Adaptation in Natural and Artificial Systems}
(1975) gave us the genetic algorithm: a population of candidate
solutions, a fitness function, selection of the fittest, mutation,
recombination, repeat. The genetic algorithm was originally a tool for
studying evolution. It has since become a tool for optimisation:
routing problems, scheduling problems, neural-network architecture
search, antenna design. NASA's ST5 spacecraft flew an antenna whose
shape was found by a genetic algorithm and would not have been designed
by a human engineer; it looks like a bent paperclip.

There is something epistemically awkward about the field of
evolutionary computation. The same machinery --- variation, selection,
inheritance --- is used both to study how natural evolution produces
life-like structure and to find good neural-network architectures, and
the second use is enormously more lucrative than the first. A genetic
algorithm that gets stuck on a local optimum, in the optimisation
literature, is failing at its job. A genetic algorithm that converges,
in the artificial-life literature, is failing at its job in a different
sense: it has stopped producing novelty.

Evolutionary computation makes its appearance in this book primarily
as a set of variation operators. Crossover, mutation, tournament
selection --- these are concrete realisations of the second slot of
the lens, and a reader who has ever written a GA already has hands-on
experience with the questions Chapter 5 makes precise.

\section{Four questions the rest of the book takes seriously}

Across the five systems we have toured, four questions keep recurring.

\textit{What makes any of this seem alive?} The Lenia blob, the Boids
flock, the Tierra parasite --- they each provoke the same uncomfortable
intuition that \emph{something} is going on at a level above the rules.
None of them is alive, and we are not going to argue otherwise. But the
intuition is not nothing. Chapter 2 takes the philosophical question of
life seriously, surveys the live frames from Aristotle through Friston,
and previews how the book will use them.

\textit{What is emergence?} The gliders are emergent in some sense; the
flocking is emergent in some sense; the parasite ecology is emergent in
some sense. Are these the same sense? Chapter 3 surveys the frames ---
weak versus strong emergence, supervenience, predictive irreducibility
--- and lays out the disagreements without resolving them.

\textit{What is novelty?} Tierra produces novelty for a while and then
stops. The Game of Life produces novelty when you change the initial
state, but a single run reaches stable patterns and is done. Chapter 3
tackles this too: novelty as the production of unseen states, and the
detection problem that makes it hard to study.

\textit{What would open-endedness mean?} The thing natural evolution
seems to do that none of our artificial systems quite manages:
indefinite novelty production. Chapter 8 surveys current proposals for
measuring it, including the $\Omega$ metric, and Chapter 9 argues that
the problem is still meaningfully open.

These are not questions the book answers. They are the questions the
book gives you tools to \emph{read}. After Chapter 4 you will be able
to take any artificial-life system --- including ones that have not
been built yet --- and write down its substrate, its variation
operator, its viability filter, its interaction topology, and its
observer; and then you will be able to read each of the four questions
through that decomposition and watch the disagreements organise
themselves in a way that is much easier to think about than the bare
arguments.

That is the deal the book is offering. The next two chapters are about
the questions. Chapter 4 is the lens. The rest is what the lens lets
you do.

\section*{To think about}

\begin{enumerate}
  \item The Game of Life and Lenia are sometimes called discrete and
    continuous versions of ``the same idea.'' What does \emph{the same}
    do here? Where does the analogy strain?
  \item Reynolds' Boids has three rules. Try to write down a fourth
    that might plausibly be in the algorithm but is not. How does it
    change the flocking behaviour? (If you have access to a Boids
    notebook, try it.)
  \item Tierra's seed program was 80 instructions, hand-written. The
    first parasites that emerged were 45 instructions. Why is
    ``shorter than the seed'' a \emph{theoretically} unsurprising
    direction for evolution to go, even though Ray did not anticipate
    it?
  \item Evolutionary computation uses the same machinery for studying
    life and for optimising neural networks. Pick a third domain ---
    protein design, route planning, chemistry --- and describe how a
    GA is being used there. Is it a study of life or an exercise in
    optimisation? Could it be both?
\end{enumerate}
